\chapter{État de l'art et positionnement}
\label{chap:etatart}

\section{La maintenance conditionnelle et l'architecture OSA-CBM}
La maintenance conditionnelle s'appuie sur une chaîne de traitement standardisée par la
norme \textbf{ISO~13374}, connue sous le nom d'architecture \textbf{OSA-CBM} (\emph{Open
System Architecture for Condition-Based Maintenance}). Cette architecture décompose le
processus en couches fonctionnelles successives :
\begin{enumerate}[leftmargin=1.6em]
  \item \textbf{Data Acquisition} (DA) : acquisition des signaux bruts.
  \item \textbf{Data Manipulation} (DM) : traitement du signal, extraction d'indicateurs.
  \item \textbf{State Detection} (SD) : comparaison à des seuils, détection d'anomalies.
  \item \textbf{Health Assessment} (HA) : évaluation de l'état de santé global.
  \item \textbf{Prognostics} (PA) : pronostic, estimation de la durée de vie restante.
  \item \textbf{Advisory Generation} (AG) : génération de recommandations et d'actions.
\end{enumerate}
PrognoSense reprend fidèlement cette décomposition : l'ingestion (DA), l'extraction
d'indicateurs (DM), la détection d'anomalies et le seuillage (SD), l'indice de santé (HA),
le pronostic RUL (PA) et la génération d'ordres de travail (AG).

\section{Techniques d'analyse vibratoire}
L'analyse vibratoire est l'outil de diagnostic privilégié des machines tournantes car
chaque mode de défaillance produit une \emph{signature} reconnaissable :
\begin{itemize}[leftmargin=1.5em]
  \item \textbf{Analyse temporelle} : des descripteurs statistiques (RMS, kurtosis, facteur
        de crête) caractérisent l'énergie et le caractère impulsif du signal. Le kurtosis
        et le facteur de crête sont particulièrement sensibles aux chocs répétés des défauts
        de roulement naissants.
  \item \textbf{Analyse fréquentielle} : la transformée de Fourier (FFT) révèle les
        composantes spectrales. Le balourd se manifeste à la fréquence de rotation
        ($1\times$), le désalignement à $2\times$, le jeu mécanique par une multiplicité
        d'harmoniques.
  \item \textbf{Fréquences caractéristiques de roulement} : BPFO (bague externe), BPFI
        (bague interne), BSF (billes), FTF (cage), calculées à partir de la géométrie du
        roulement et de la vitesse de rotation.
  \item \textbf{Analyse d'enveloppe} : la démodulation par transformée de Hilbert d'une
        bande de résonance haute fréquence fait ressortir les défauts de roulement masqués
        dans le spectre brut.
\end{itemize}

\section{Apprentissage automatique pour le diagnostic et le pronostic}
Au-delà du diagnostic expert manuel, l'apprentissage automatique permet d'automatiser la
reconnaissance des défauts (classification supervisée) et l'estimation de la durée de vie
restante (régression). Les approches non supervisées (auto-encodeurs, méthodes de détection
d'anomalies) complètent ce dispositif en repérant les comportements \emph{inconnus}, non
présents dans les données d'entraînement --- un atout décisif, puisqu'une partie importante
des défaillances industrielles relèvent de modes rares ou imprévus.

\section{Normes applicables}
Le domaine est fortement encadré par des normes que PrognoSense intègre explicitement :
\begin{itemize}[leftmargin=1.5em]
  \item \textbf{ISO~10816 / 20816} : évaluation de la sévérité vibratoire par la vitesse
        efficace (mm/s) et classement en zones A/B/C/D selon la classe de machine.
  \item \textbf{ISO~13374} : architecture de traitement OSA-CBM (ci-dessus).
  \item \textbf{ISO~13373} : surveillance et diagnostic vibratoire.
  \item \textbf{ISO~13381} : pronostic.
  \item \textbf{ISO~18436} : qualification des analystes vibratoires (catégories I à IV).
\end{itemize}

\section{Solutions du marché et positionnement de PrognoSense}
Le marché de la maintenance prédictive est occupé par des acteurs établis : SKF
(Enlight, @ptitude), Emerson (AMS, Plantweb), GE Bently Nevada (System~1), Brüel \& Kjær,
ainsi que des acteurs plus récents centrés sur l'IA et les capteurs sans fil (Augury,
AspenTech~Mtell, Siemens~Senseye, Schaeffler~OPTIME). Ces solutions sont puissantes mais
présentent des limites récurrentes : matériel propriétaire, coût de licence élevé,
fermeture (peu d'interopérabilité), et fréquemment un caractère de \og boîte noire \fg{}
peu explicable.

PrognoSense ne vise pas à rivaliser en largeur fonctionnelle avec ces suites, mais à
occuper un positionnement différenciant, articulé autour de cinq axes :
\begin{enumerate}[leftmargin=1.6em]
  \item \textbf{Ouverture} : connectable à n'importe quel capteur via les protocoles
        standards (OPC-UA, MQTT), sans verrouillage propriétaire.
  \item \textbf{Diagnostic normalisé} : sévérité exprimée en zones ISO~10816, défendable
        par un ingénieur de fiabilité.
  \item \textbf{Explicabilité et confiance} : importance des indicateurs, calibration,
        journal d'audit, taux de fausses alarmes mesuré.
  \item \textbf{Copilot ancré} : assistant en langage naturel dont les réponses citent
        les sources normatives (RAG).
  \item \textbf{Boucle fermée} : de l'alerte à l'ordre de travail transmis à la GMAO.
\end{enumerate}
Ce positionnement guide l'ensemble des choix de conception détaillés dans les chapitres
suivants.
